	\section{Рабочий проект}
Рабочий проект является одним из ключевых разделов Выпускной квалификационной работы по теме Бизнес-Проект "Закрытый корпоративный мессенджер. Разработка клиентской части". Поскольку он описывает непосредственную реализацию разработанной системы. Переход от теоретических обоснований и проектных решений к конкретному воплощению.  
\subsection{Спецификация модулей}
Закрытый корпоративный мессенджер спроектирован как модульная система, обеспечивающая гибкость разработки и масштабируемость. Система разделена на логические модули, каждый из которых отвечает за определенный функционал и может быть развит независимо. Ниже представлено описание ключевых модулей.
\subsubsection{Модуль регистрации}
Модуль регистрации обеспечивает создание новых учетных записей в мессенджере, реализуя полный цикл от ввода данных пользователем до обработки на сервере. Его архитектура включает три компонента: register.html (структура интерфейса), register.css (визуальное оформление) и register.js (бизнес-логика). Ключевая задача модуля - безопасная обработка регистрационных данных с минимизацией ошибок пользователя.

HTML-основа (register.html) построена вокруг центрального контейнера формы с семантическими элементами. Форма содержит два обязательных поля ввода: для логина и пароля, кнопку отправки данных, а также ссылку для перехода к странице входа. Особое внимание уделено контейнеру alerts-container, который динамически отображает уведомления поверх основного контента без перезагрузки страницы. Скрытых полей или сложных структур нет - дизайн сознательно минимизирован для снижения когнитивной нагрузки.

Визуальный слой (register.css) использует единую стилистику с модулем авторизации, сохраняя градиентный фон и полупрозрачный белый контейнер формы. Акценты расставлены через:
\begin{itemize}
	\item контрастную обводку поля ввода при ошибках (красный цвет);
	\item плавные переходы состояний элементов;
	\item анимацию появления уведомлений;
	\item адаптивную верстку для мобильных устройств;
	\item интерактивные эффекты для ссылок.
\end{itemize}
Логика обработки (register.js) начинается с инициализации после полной загрузки DOM. Ядро модуля - обработчик события submit формы, который выполняет многоуровневую валидацию. Клиентская валидация логина проверяющее:
\begin{itemize}
	\item допустимые символы (латиница, кириллица, цифры, дефисы);
	\item длину строки (3-15 символов);
	\item отсутствие спецсимволов.
\end{itemize}
Проверка пароля на минимальную длину (6 символов) и наличия значения. Визуальная индикация ошибок:
\begin{itemize}
	\item добавление класса .input-error для проблемных полей;
	\item динамическое изменение фона поля;
	\item автоматический фокус на некорректном поле.
\end{itemize}

Система уведомлений реализована через функцию showAlert, которая создает тосты трех типов (ошибка, успех, информация) с иконками состояния. Уведомления появляются в верхней части экрана с анимацией slide-in и автоматически исчезают через 4 секунды. Особенность реализации - каскадное удаление уведомлений через очередь анимаций.

Модуль интегрирован с системой аутентификации через ссылку перехода к странице входа, сохраняя единую стилистику интерфейса. Все взаимодействия спроектированы с учетом важных элементов (кнопка отправки) имеют увеличенные размеры и контрастное оформление.
\subsubsection{Модуль авторизации}
Модуль авторизации является критически важным компонентом системы, обеспечивающим безопасный доступ пользователей к мессенджеру. Его реализация включает три ключевых файла: login.html (интерфейс), login.css (стилизация) и login.js (логика работы). Основная задача модуля - аутентификация пользователей через проверку учетных данных с последующим перенаправлением в основной интерфейс мессенджера.

HTML-структура (login.html) представляет собой минималистичную форму входа с двумя полями ввода: для имени пользователя и пароля. Элемент alerts-container динамически отображает системные уведомления об ошибках валидации или проблемах аутентификации. Ключевой особенностью является скрытое поле user-id, которое после успешной аутентификации заполняется серверными данными для последующей идентификации пользователя в системе. Форма использует семантическую разметку с четкой иерархией элементов, что улучшает доступность.

Стилевое оформление (login.css) построено на мягкой пастельной цветовой палитре с градиентным фоном, создающим визуально приятную атмосферу. Контейнер формы имеет полупрозрачный белый фон со скругленными углами и выраженной теневой подсветкой, что обеспечивает "парящий" эффект. Интерактивные элементы реализуют сложные анимации: кнопка входа имеет динамичный градиентный фон с эффектом "блеска" при наведении, реализованным через псевдоэлемент ::after с анимацией shine. Поля ввода анимируют плейсхолдеры - при фокусировке текст плавно смещается вверх с одновременным затуханием. Для мобильных устройств применена адаптивная верстка с перестроением элементов и оптимизацией отступов.

Логика работы (login.js) начинается с обработки события DOMContentLoaded, гарантирующей инициализацию скриптов после полной загрузки DOM. Ядром модуля является обработчик события submit формы, который выполняет:
\begin{enumerate}
	\item Валидацию имени пользователя через регулярное выражение [a-zA-Zа-яА-Я0-9\_-]{3,15}, запрещающее спецсимволы и ограничивающее длину логина.
	\item Формирование тела запроса в формате application/x-www-form-urlencoded.
	\item Отправку данных на эндпоинт /login через Fetch API.
	\item Обработку ответа сервера с сохранением сессионных данных. session id и session key в sessionStorage. Имя пользователя для отображения в интерфейсе.
	\item Перенаправление на главную страницу при успешной аутентификации.
	\item Визуализацию ошибок через функцию showAlert.
\end{enumerate}
Система уведомлений реализована через динамическое создание DOM-элементов. Функция showAlert генерирует тосты с иконками состояния, которые автоматически исчезают через 5 секунд с применением CSS-анимации fadeIn. Особое внимание уделено обработке сетевых ошибок - при обрыве соединения пользователь получает понятное сообщение об ошибке соединения.
Модуль полностью адаптирован под мобильные устройства с применением медиазапросов, оптимизирующих расположение элементов на экранах меньше 480px. Все элементы интерфейса имеют увеличенные области клика для удобства touch-интерфейса.
\subsubsection{Ядро чата. Управление чатами}
Управление чатами обеспечивает организацию коммуникационных пространств через централизованную систему хранения состояния в sessionStorage, автоматически восстанавливающую последний активный чат при загрузке приложения. Групповые чаты реализуют иерархическую модель ролей (владелец/администратор/участник), динамически адаптируя интерфейс под права пользователя - владельцы управляют составом и параметрами групп, администраторы регулируют участников, рядовые члены ограничены отправкой сообщений. Личные чаты создаются автоматически при первом обращении через поиск пользователей, с интеллектуальным кэшированием новых диалогов для мгновенного отклика.

Переключение между чатами инициирует каскад действий:
\begin{itemize}
	\item сброс текущих сообщений;
	\item установку нового контекста;
	\item загрузку истории и обновление интерфейса;
	\item с автоматической проверкой прав доступа к групповым чатам. 
\end{itemize}

Управление участниками включает динамическую подгрузку членов групп с дебаунсингом запросов и контекстным меню для администраторов, реализующим изменение ролей и исключение пользователей. Модальные окна создания/модификации чатов выполняют валидацию ввода и обрабатывают исключительные сценарии - при потере доступа система автоматически переключает пользователя в общий чат, предотвращая зависание в несуществующем контексте.
\subsubsection{Ядро чата. Работа с сообщениями}
Функции обработки сообщений реализуют полный жизненный цикл контента. Система отправки через sendMessage() поддерживает текстовые сообщения и файловые вложения (изображения, документы) с предварительным отображением временных сообщений. Для личных чатов применяется AES-шифрование через CryptoJS с использованием session\-key. Загрузка истории реализована в loadMessages()/loadPrivateMessages() с инкрементальной подгрузкой по временным меткам. Редактирование выполняется через контекстное меню с преобразованием текста в интерактивное поле ввода. Удаление сообщений требует подтверждения в модальном окне. Отображение в displayMessages() включает:
\begin{itemize}
	\item автоматическую прокрутку к новым сообщениям;
	\item стилизацию системных уведомлений;
	\item визуализацию вложений с предпросмотром изображений;
	\item подсветку измененных сообщений.
\end{itemize}
\subsubsection{Ядро чата. UI взаимодействия}
Интерфейсные функции обеспечивают интуитивное взаимодействие. Эмодзи-пикер (initEmojiPicker()) реализует сетку смайлов с анимацией выбора. Контекстное меню сообщений (showContextMenu()) динамически формирует опции:
\begin{itemize}
	\item редактирование/удаление текста;
	\item просмотр и скачивание вложений;
	\item массовая загрузка файлов;
	\item система уведомлений.
\end{itemize}
\subsubsection{Интерфейс}
Основные стили мессенджера построены на адаптивной grid-сетке с использованием флексбокс-контейнеров, обеспечивающей стабильную структурную целостность на всех типах устройств. Цветовая схема сознательно выбрана в пастельной гамме с доминирующими оттенками персикового и розового для фоновых элементов, контрастирующими с глубоким шоколадным цветом текста, что создает визуально комфортную рабочую среду без излишней яркости. Анимационная система реализует сложные взаимодействия через плавные CSS-переходы для интерактивных элементов, трансформации с 3D-эффектами для кнопок и модальных окон, кастомные keyframe-анимации для входящих сообщений (messageSlide) и подсветки измененного контента (highlight), а также динамические градиенты с мягкими цветовыми переходами. Типографика использует безопасный стек шрифтов (Arial, sans-serif) с адаптивными размерами текста и четкой визуальной иерархией заголовков, где эффекты глубины достигаются через многоуровневые тени (box-shadow) с переменной прозрачностью, создающие иллюзию "парящих" элементов интерфейса над фоном.

Специализированные компоненты разработаны с учетом функциональных требований и эстетического единства. Стили сообщений реализуют четкую дифференциацию типов контента: пользовательские сообщения оформлены как бирюзовые блоки с треугольным указателем и теневой подсветкой, системные уведомления представлены центрированными серыми карточками с курсивным текстом и полупрозрачным фоном, а собственные сообщения пользователя выделены контрастным синим цветом с правым выравниванием и специальным градиентом. Дизайн сайдбаров сочетает практическую функциональность с визуальной привлекательностью через градиентные фоны, усиленные мягкими теневыми акцентами, где интерактивные элементы обеспечивают тактильную обратную связь при взаимодействии. Модальные окна реализованы как "парящие" карточки со сложными анимациями появления (modalAppear), включающими одновременное масштабирование и смещение, с контрастными акцентными элементами и продуманной типографической иерархией. Интерактивные элементы используют составные эффекты: кнопки сочетают градиентные фоны, трансформации при наведении (scale, translateY), динамические тени и псевдоэлементы для создания эффектов свечения ("shine"), а поля ввода анимируют плейсхолдеры при фокусировке с одновременным смещением и затуханием.

Адаптивность системы реализована через mobile-first подход с тремя ключевыми контрольными точками (1200px, 768px, 480px), обеспечивающий оптимальное отображение на всех классах устройств. Для мобильных экранов применены радикальные трансформации: навигационная панель перестраивается в компактный горизонтальный контейнер с иконографическими кнопками, сайдбары занимают 85 процентов ширины экрана с выездной анимацией, а компоновка элементов формы ввода сообщения адаптируется под вертикальную схему с сохранением удобства взаимодействия. Все интерактивные элементы имеют увеличенные области касания (min-width: 44px) с переработанной навигацией, которая заменяет текстовые метки на интуитивные иконки на малых экранах. Медиа-запросы оптимизируют не только размеры, но и взаимодействие: на устройствах с шириной менее 768px автоматически отключаются сложные фоновые эффекты для повышения производительности, корректируется высота чат-контейнера для освобождения пространства под виртуальную клавиатуру, и активируется специальная версия контекстного меню, активируемого долгим нажатием. Техники адаптивной типографики включают плавное уменьшение размеров шрифта, изменение межбуквенных интервалов и динамическую регулировку высоты строки для сохранения читаемости контента на экранах любой диагонали.

\subsubsection{Вспомогательные модули}
Обработка ошибок реализована через кастомные страницы 403, 404 и 500, оформленные в едином визуальном стиле с использованием error.css. Все страницы ошибок сохраняют фирменную пастельную цветовую гамму и типографику основного интерфейса, обеспечивая последовательность пользовательского опыта даже в исключительных ситуациях. Каждая страница содержит понятное описание проблемы (например, "Доступ запрещен" для 403 или "Страница не найдена" для 404), сопровождаемое релевантной иконкой или анимацией, и главное - навигационную ссылку для возврата на домашнюю страницу, что позволяет пользователю быстро восстановить работоспособность сессии без использования браузерной кнопки "Назад".

Управление файлами включает комплексную систему обработки вложений, начиная с прелоадера изображений, который отображает плейсхолдеры во время загрузки контента. Система предпросмотра вложений реализована через модальные окна, позволяющие просматривать изображения в полноразмерном режиме и скачивать документы без перехода в чат. Индикатор выбранных файлов динамически отображает количество и типы загружаемых вложений прямо над полем ввода сообщения, с функцией очистки всего выбора одним кликом. Для файлового меню реализована интеллектуальная группировка - изображения отображаются в виде миниатюр с возможностью масштабирования, а документы представлены иконками с четким указанием формата и размера.

Системные компоненты обеспечивают целостность интерфейса через библиотеку иконок, созданных средствами CSS и SVG для сохранения четкости на любых разрешениях экрана. Шаблоны модальных окон используют единую архитектуру с параметризуемыми заголовками, контентом и действиями, что позволяет быстро создавать диалоги подтверждения, формы ввода и информационные панели. Все шаблоны реализуют сложную анимационную логику с плавным появлением, затемнением фона и адаптивным позиционированием, включая автоматическую подстройку под размеры мобильных устройств с выносом критических элементов в зону удобного доступа.

\subsection{Тестирование интерфейса и функций мессенджера}
Тестирование интерфейса и функций веб-приложения "Закрытый корпоративный мессенджер" - это критически важный этап разработки, который обеспечивает стабильность, надежность и удобство использования веб-приложения. 
\subsubsection{Регистрация и авторизация пользователя}
Необходимо убедиться, что процессы регистрации нового пользователя и авторизации работают корректно в соответствии с требованиями, обеспечивая надежный доступ к приложению. Программа должна обрабатывать все возможные ситуации и события.

Описание: пользователь открывает страницу регистрации /register, вводит логин и пароль и нажимает кнопку "Зарегистрироваться".


Ожидаемый результат: при успешной регистрации с корректными данными система перенаправляет на страницу авторизации. При вводе некорректных данных система уведомляет пользователя об этом.

Пользователь нажал на кнопку "Зарегистрироваться"  без введенных данных (рис. \ref{fig:ui-registration1}).
\newpage
\begin{figure}[ht]
	\centering
	\includegraphics[width=0.4\linewidth]{"images/Заполните поля регистрация"}
	\caption{Уведомление для пользователя о вводе данных в поля}
	\label{fig:ui-registration1}
\end{figure}

Пользователь ввел логин, который уже был зарегистрирован (рис. \ref{fig:ui-registration2}).
\begin{figure}[ht]
	\centering
	\includegraphics[width=0.8\linewidth]{"images/Имя пользователя"}
	\caption{Уведомление для пользователя о вводе логина, который уже зарегистрирован}
	\label{fig:ui-registration2}
	\end{figure}

Проверка пароля. Пользователь ввел менее 6 символов в пароль (рис. \ref{fig:ui-registration3}).
\begin{figure}[ht]
	\centering
	\includegraphics[width=0.8\linewidth]{"images/Пароль"}
	\caption{Уведомление для пользователя о вводе пароля менее 6 символов}
	\label{fig:ui-registration3}
\end{figure}

Успешная регистрация пользователя и переадресация на страницу Входа /login (рис. \ref{fig:ui-registration4}).
\newpage
\begin{figure}[ht]
	\centering
	\includegraphics[width=0.8\linewidth]{"images/Регистрация"}
	\caption{Уведомление пользователю об успешной регистрации}
	\label{fig:ui-registration4}
\end{figure}

Пользователь ввел некорректный логин при регистрации (рис. \ref{fig:ui-registration5}).
\begin{figure}[ht]
	\centering
	\includegraphics[width=0.8\linewidth]{"images/Пароль15"}
	\caption{Уведомление о некорректном логине при регистрации}
	\label{fig:ui-registration5}
\end{figure}


Описание: пользователь после регистрации попадает на страницу авторизации (вход) /login. Пользователь вводит свои данные, которые вводил при регистрации (логин, пароль).


Ожидаемый результат: при успешной авторизации и вводе корректных данных пользователя, открывается главная страница веб-приложения. При вводе некорректных данных система уведомляет пользователя об этом.
Пользователь ввел некорректный или неверный пароль в поле ввода (рис. \ref{fig:ui-auth2}).
\begin{figure}[ht]
	\centering
	\includegraphics[width=0.8\linewidth]{"images/Неверный пароль"}
	\caption{Уведомление пользователю о неверном пароле}
	\label{fig:ui-auth2}
\end{figure}


Пользователь ввел логин незарегистрированного пользователя (рис. \ref{fig:ui-auth3}). 
\begin{figure}[ht]
	\centering
	\includegraphics[width=0.8\linewidth]{"images/Пользователь не найден"}
	\caption{Уведомление пользователю о незарегистрированном логине}
	\label{fig:ui-auth3}
\end{figure}


Успешный вход переадресовывает на главную страницу веб-приложения. В правом верхнем углу главной страницы указывается логин авторизованного пользователя (рис. \ref{fig:ui-auth4}).
\begin{figure}[ht]
	\centering
	\includegraphics[width=0.8\linewidth]{"images/Логин пользователя"}
	\caption{Логин успешно авторизованного пользователя на главной странице}
	\label{fig:ui-auth4}
\end{figure}
\begin{xltabular}{\textwidth}{|X|l|X|}
	\caption{Результаты тестирования регистрации и авторизации\label{table}}\\ \hline
	\centrow Сценарий & \centrow Статус &  \centrow Описание \\ \hline
			\endfirsthead
	\continuecaption{Продолжение таблицы \ref{table}}
		\centrow Сценарий & \centrow Статус &  \centrow Описание \\ \hline
		\finishhead
	Проверка на заполняемость полей & Успешно & Система уведомляет пользователя "Пожалуйста, заполните это поле." \\ \hline
	Проверка на существующий логин & Успешно & Система уведомляет пользователя "Пользователь с таким именем уже существует" \\ \hline
	Проверка пароля & Успешно & Система уведомляет пользователя "Пароль должен быть не менее 6 символов" \\ \hline
	Проверка на корректный логин & Успешно & Система уведомляет о некорректном вводе логина \\ \hline
	Уведомление о регистрации & Успешно & Система уведомляет пользователя о "Регистрация прошла успешно! Перенаправление..." и перенаправляет на авторизацию \\ \hline 
	Проверка пароля & Успешно & Система уведомляет "Неверный пароль" \\ \hline
	Проверка логина & Успешно & Система уведомляет "Пользователь не найден" \\ \hline
	Успешная авторизация & Успешно & Перенаправляет на главную страницу, выводит "Вы вошли: логин"  
\end{xltabular}

\subsubsection{Главная страница веб-приложения. Тестирование чат-блока}
В этом подразделе описываются методы и результаты тестирования чат-блока на главной странице веб-приложения, включая отправку и получение сообщений, отображение истории переписки, работу с файлами  и интерактивные элементы управления, для подтверждения его полной функциональности и удобства использования.

Описание: пользователь вводит большого размера текст в поле ввода отправки сообщений.


Ожидаемый результат: поля ввода имеет первоначальный фиксированный размер, но при вводе текста с большим количеством символов поле растягивается до заранее ограниченного размера. Появляется индикатор прокрутки, который позволяет перемещаться по тексту от начала до конца (скроллить).

Поле ввода сообщений с фиксированным размером (рис. \ref{fig:ui-chat1}).
\begin{figure}[ht]
	\centering
	\includegraphics[width=0.8\linewidth]{"images/Фикс чат"}
	\caption{Фиксированный размер поля ввода сообщения}
	\label{fig:ui-chat1}
\end{figure}

Поле ввода сообщений растянутое до ограниченных размеров, появляется бегунок скролла (рис. \ref{fig:ui-chat2}).
\begin{figure}[ht]
	\centering
	\includegraphics[width=0.8\linewidth]{"images/Большой чат"}
	\caption{Растянутое поле ввода сообщений с полосой прокрутки}
	\label{fig:ui-chat2}
\end{figure}

Описание: пользователь отправил сообщение, но хочет его отредактировать.


Ожидаемый результат: пользователь нажимает правой кнопкой по своему отправленному сообщению, после чего редактирует его (рис. \ref{fig:ui-edit1}) и сохраняет изменения. Новое отредактированное сообщение верно отображается в чате. 

\begin{figure}[ht]
	\centering
	\includegraphics[width=0.8\linewidth]{"images/Редактирование"}
	\caption{Редактирование отправленного сообщения}
	\label{fig:ui-edit1}
\end{figure}

Описание: пользователь прикрепляет желаемый файл для отправки и отправляет его.


Ожидаемый результат: пользователь нажимает на кнопку "Файл" (обозначенный как скрепка), прикрепляет желаемый файл, который имеет предпросмотр. При нажатии кнопки "Отправить" прикрепленный файл отправляет в чат (рис. \ref{fig:ui-file}).
\newpage
\begin{figure}[ht]
	\centering
	\includegraphics[width=0.8\linewidth]{"images/Файл"}
	\caption{Предпросмотр прикрепленного файла}
	\label{fig:ui-file}
\end{figure}

Описание: пользователь хочет посмотреть вложения отправленного сообщения с прикрепленным файлом.


Ожидаемый результат: пользователь нажимает правую кнопку по отправленному сообщению в чате, далее выбирает "Просмотреть вложения". Система открывает отдельное модальное окно с просмотром вложений и возможностью скачать вложения (рис. \ref{fig:ui-file1}). 

\begin{figure}[ht]
	\centering
	\includegraphics[width=0.8\linewidth]{"images/Вложения"}
	\caption{Просмотр вложений прикрепленный к отправленному сообщению в чате}
	\label{fig:ui-file1}
\end{figure}

\begin{xltabular}{\textwidth}{|X|l|X|}
	\caption{Результаты тестирования отправки сообщений}\\ \hline
	\centrow Сценарий & \centrow Статус &  \centrow Описание \\ \hline
	Большое количество символов в тексте при отправке & Успешно & Поля ввода сообщений растягивается, появляется ползунок скролла \\ \hline
	Редактирование отправленного сообщения & Успешно & Сообщение редактируется, подсвечивается в чате, моментально обновляется \\ \hline
	Прикрепление файла & Успешно &  Прикрепленный файл предпросматривается, отображается количество файлов и отменя отправки \\ \hline
	Просмотр вложений & Успешно & Отправленные файлы просматриваются, скачиваются \\ \hline
	Удаление сообщений & Успешно & Сообщения удаляются, моментально исчезают из чата \\ \hline 
\end{xltabular}


\subsubsection{Главная страница. Сайдборды, боковые меню} 

В данном подразделе рассматривается тестирование функциональности боковых меню (сайдбордов) - левого с группами и личными чатами и правого с участниками, включая их открытие и закрытие, переключение между состояниями и корректное отображение содержимого на всех типах устройств.

Описание: пользователь открыл боковые меню, левое с группами и личными чатами, правое с участниками активного чата.


Ожидаемый результат: пользователь увидел полную информацию о групповых и личных чата, об участниках активного чата. Независимо от размера экрана пользователя, информация предоставляется полностью, без искажения (рис. \ref{fig:ui-sidebar}).

\begin{figure}[ht]
	\centering
	\includegraphics[width=0.8\linewidth]{"images/Боковые ПК"}
	\caption{Боковые меню на компьютере}
	\label{fig:ui-sidebar}
\end{figure}

Левое боковое меню на мобильном устройстве (рис. \ref{fig:ui-sidebar1}).

\begin{figure}[ht]
	\centering
	\includegraphics[width=0.3\linewidth]{"images/Левое тф"}
	\caption{Левое боковое меню с группами и личными чата на мобильном устройстве}
	\label{fig:ui-sidebar1}
\end{figure}

\newpage

Правое боковое меню на мобильном устройстве (рис. \ref{fig:ui-sidebar2}). 

\begin{figure}[ht]
	\centering
	\includegraphics[width=0.4\linewidth]{"images/Правое тф"}
	\caption{Правое боковое меню с участниками активного чата на мобильном устройстве}
	\label{fig:ui-sidebar2}
\end{figure}

Описание: пользователь с ролью "Владелец" (создатель чата) хочет изменить роль другого участника чата.


Ожидаемые результат: пользователь изменяет роль другого участника чата (рис. \ref{fig:ui-role1}). Участник получает новую роль "Администратор". Владелец чата получает уведомление об успешном изменении роли, отдельно сообщение об изменении в групповом чате.
\newpage
\begin{figure}[ht]
	\centering
	\includegraphics[width=0.8\linewidth]{"images/Роль"}
	\caption{Изменение роли у другого участника чата}
	\label{fig:ui-role1}
\end{figure}


Диалоговое окно выбора роли для другого участника чата (рис. \ref{fig:ui-role2}).

\begin{figure}[ht]
	\centering
	\includegraphics[width=0.5\linewidth]{"images/Диалог роли"}
	\caption{Меню выбора роли для участника чата}
	\label{fig:ui-role2}
\end{figure}

Уведомление об успешном изменении роли. Сообщение в групповом чате об изменении роли для определенного участника (рис. \ref{fig:ui-role3}).
\newpage
\begin{figure}[ht]
	\centering
	\includegraphics[width=0.8\linewidth]{"images/Уведомление роли"}
	\caption{Уведомление об изменении роли и сообщение об этом в групповой чат}
	\label{fig:ui-role3}
\end{figure}


Описание: пользователь желает создать новый групповой чат.


Ожидаемый результат: пользователь нажимает отведенную кнопку для создания группового чата "Новая группа". Система выводит отдельное модальное окно "Создать новую группу", вместе с полем ввода "Название группы" и двумя кнопками "Отмена" и "Создать". Группа создается, открывается чат этой группы, в левом боковом меню появляется название этой группы, в правом боковом меню участник (владелец этой группы).

Кнопка создания группы "Новая группа" в левом боковом меню (рис. \ref{fig:ui-group1}).
\newpage
\begin{figure}[ht]
	\centering
	\includegraphics[width=0.5\linewidth]{"images/Кнопка группа"}
	\caption{Кнпока Новая группа в левом боковом меню}
	\label{fig:ui-group1}
\end{figure}


Модальное окно создания новой группы с полем ввода названия (рис. \ref{fig:ui-group2}).

\begin{figure}[ht]
	\centering
	\includegraphics[width=0.5\linewidth]{"images/Окно группы"}
	\caption{Модальное окно создания группы}
	\label{fig:ui-group2}
\end{figure}

Уведомление для пользователя о том, что группа с таким названием уже создана (рис. \ref{fig:ui-group3}).

\begin{figure}[ht]
	\centering
	\includegraphics[width=0.5\linewidth]{"images/Уведомление группы"}
	\caption{Уведомление о существующем названии группы}
	\label{fig:ui-group3}
\end{figure}

Описание: пользователь хочет начать личный чат с участником общего чата.


Ожидаемый результат: пользователь открывает правое боковое меню с участниками, нажимает левой кнопкой мыши по пользователю с кем хочет начать личный чат. Открывает личный чат и отображается в левом боковом меню (рис. \ref{fig:ui-private}).

\begin{figure}[ht]
	\centering
	\includegraphics[width=0.4\linewidth]{"images/Личный"}
	\caption{Личный чат с пользователем и его отображение в левом боковом меню.}
	\label{fig:ui-private}
\end{figure}
\newpage
\subsubsection{Результаты тестирование интерфейса и функций веб-приложения}

\begin{xltabular}{\textwidth}{|X|l|X|}
	\caption{Результаты тестирования интерфейса}\\ \hline
	\centrow Сценарий & \centrow Статус &  \centrow Описание \\ \hline
	Регистрация и авторизация & Успешно & Пользователь регистрируется и входит \\ \hline
	Отправка сообщений и их редактирование & Успешно & Сообщения отправляются и редактируются \\ \hline
	Предпросмотр файлов и скачивание вложений & Успешно & Файлы предпросматриваются и скачиваются \\ \hline
	Уведомление пользователей & Успешно & Пользователи уведомленяются о совершении действий и ошибках \\ \hline 
	Левое боковое меню & Успешно & Открывается, на мобильных устройствах показывает всю информацию \\ \hline
	Правое боковое меню & Успешно & Открывается, на мобильных устройствах показывает всю информацию \\ \hline
	Изменение роли участника & Успешно & Изменяются, приходят уведомления \\ \hline
	Создание группового чата & Успешно & Создается, информация отображается \\ \hline
	Личный чат & Успешно & Создается, отображается в меню \\ \hline
\end{xltabular}

Проведенное всестороннее тестирование интерфейса и функциональности веб-приложения "Закрытый корпоративный мессенджер" подтвердило его полную работоспособность и соответствие заявленным требованиям. В ходе тестирования были успешно проверены все ключевые сценари: начиная от базовых процессов регистрации и авторизации пользователей до сложной логики отправки и получения сообщений в реальном времени, управления группами и обработки файлов.

Тестирование выявило, что пользовательский интерфейс интуитивно понятен, визуально консистентен и обеспечивает комфортное взаимодействие на различных устройствах благодаря реализованной мобильной адаптации. Все интерактивные элементы, такие как кнопки, сайдбары и модальные окна, функционируют корректно, предоставляя пользователю ожидаемую обратную связь.
